
\section{Introduction}














Many reinforcement-learning algorithms follow the policy-iteration paradigm, which alternates between
estimating the action values of a policy
and improving the policy using those estimates.
This framework is fundamental because it can learn an optimal policy for an unknown Markov decision process (MDP) using only sampled rewards and transitions.
Classical policy iteration improves the policy once the action-value estimates have converged.
In contrast, optimistic variants improve the policy after partial evaluation.
This reduces computation as it requires fewer evaluation steps between each improvement step.
However, it complicates the analysis because

small perturbations

to the value estimates
can drastically affect the policy and alter the subsequent dynamics.


















Convergence guarantees for optimistic policy-iteration algorithms remain limited, even in the tabular setting.
Algorithms that use one-step bootstrapped estimates for evaluation, such as $Q$-learning, converge to optimality under standard asynchronous stochastic-approximation conditions.
When using multi-step bootstrapped estimates, convergence to optimality requires additional structure, such as combining several multi-step approximations~\citep{Harutyunyan2016QCorrections, Munos2016SafeLearning} or using a look-ahead mechanism~\citep{Winnicki2023OnEvaluation}.
Monte Carlo Optimistic Policy Iteration (MC-O-PI), which uses full-episode returns for evaluation, remains poorly understood. Despite its simplicity, the asymptotic behaviour of this foundational algorithm has been an open question for three decades~\citep{Sutton1999OpenLearning}.
Convergence to optimality is guaranteed when the MDP exhibits a feedforward structure \citep{Wang2022OnLearning, Lubars2021OptimisticStructure}
or when the discount factor is smaller than $1/2$ \citep{Delattre2025MarkovAlgorithm},
but, when the structure of the MDP is unknown, the only existing condition requires that all state-action pairs are updated at the same frequency \citep{Tsitsiklis2002OnIteration, Chen2018OnProblem, Liu2021OnStarts}.














Updating all pairs as frequently requires that episodes are initialised uniformly over the entire state-action space.
This condition is unrealistic in practice, especially when the state space is large.
However, once the initial state of an episode is selected, uniformly sampling the initial action is usually straightforward.
This paper proves that this weaker condition is sufficient.
Specifically, initial-visit MC-O-PI converges almost surely to the optimal action values when updates are uniform over the actions within each state.
Different states may be updated at unequal, history-dependent frequencies, provided every state accumulates infinite learning weight. We use full inertia for policy ties and a mild bound on future increases of the learning rates.
This result extends the allowable state-sampling distributions in the previous convergence guarantees
and applies in settings where uniform state sampling is not feasible but uniform action sampling is easy.






















The proof of this new theorem, stated as \autoref{thm: convergence mcopi uniform} below, requires new tools.
Previous work demonstrates convergence by studying the evolution of the value estimates. The arguments rely centrally on a commutativity property of the Bellman operator that fails when updates are not uniform over all state-action pairs.
Instead, we track the evolution of the policy and find that the mean-field dynamics of MC-O-PI generate monotonically improving policies when updates are uniform over the actions within each state.
We then prove a positive lower bound on the probability of progress during each sufficiently late interval with a fixed action-value target, stopping the estimate if it leaves a fixed bounded set. Progress means a strictly improved target or persistence of the current target. There are only finitely many targets, so repeated opportunities and stability imply that the target changes only finitely many times. The value estimate consequently converges to the optimal action values, and every sufficiently late greedy policy is optimal.

The argument develops the combined stability--ODE strategy of \citet[Section~5.4]{Kushner2003StochasticApplications} for policy-region entries whose margins may approach zero. Favorable updates create small positive gaps, and learning-rate regularity controls subsequent noise relative to those gaps. This removes the need to assume repeated entry into a fixed interior subset of a policy region. The precise local comparison with lock-in estimates is discussed before the stochastic proof.

This paper is organised as follows. Section~2 introduces finite Markov decision
processes and formulates MC-O-PI as a stochastic difference inclusion. Section~3 proves the
main convergence theorem, first for the mean-field dynamics and then for the stochastic
algorithm. Section~4 discusses the practical implications of action-wise uniform updates
and possible extensions to related optimistic methods.

Unless stated otherwise, operations and relations between functions are elementwise. For instance, given functions $f, g : \mathcal X \to \mathcal Y$ we write
\begin{align*}
&(f + g)(x) = f(x) + g(x), \qquad \forall \, x \in \xset ,
    \\
&f \le g \iff f(x) \le g(x), \qquad \forall \, x \in \xset .
\end{align*}
Moreover, we compute norms as $\|f\| := \sup_{x \in \xset} |f(x)|$
and adopt the convention $\inf \varnothing = \infty$.

























































































